\section{Related Work}
\label{sec:related}

\paragraph{Unsupervised and goal-conditioned RL.}
URL methods learn reusable control objects in a reward-free or reward-agnostic phase and reuse them for downstream goals.
Goal-conditioned methods learn values or policies indexed by target states~\citep{eysenbach-iclr21a,ma-neurips22a}, while contrastive and metric-based methods learn state-goal scores or representations that distinguish future states from negatives~\citep{eysenbach-neurips22a,wang-icml23a,park-iclr24a}.
Successor representations, successor features, and forward--backward representations estimate future-occupancy objects that can be combined with downstream rewards~\citep{dayan-nc93a,barreto-neurips17a,touati-neurips21a,touati-iclr23a}.
Proto-value and proto-successor methods instead use learned bases or spectral structure to represent behavior or value spaces~\citep{mahadevan-icml05a,mahadevan-jmlr07a,farebrother-iclr23a,agarwal-arxiv25a}.
Recent work has argued that many URL objectives can be viewed as tractable approximations to successor measures or future-reachability objects~\citep{agarwal-rlj25a}.
Our work uses this view diagnostically: we ask whether these learned objects preserve local interaction modes that change successor structure, rather than proposing a new URL objective.

\paragraph{Structured and controllable representations.}
Classical factored MDP work exploits state structure to make planning and learning more efficient~\citep{boutilier-ijcai95a,kearns-ijcai99a,guestrin-jair03a}.
More recent controllable-representation methods aim to preserve action-relevant state information, often through inverse or multi-step inverse dynamics objectives~\citep{efroni-iclr22a,lamb-arxiv22a,levine-arxiv24a,rudolph-rlc24a}.
These works are closely related to our interaction-mode view: the relevant abstraction should retain variables that change how actions affect future states.
Our contribution is different in scope.
We do not learn a new abstraction or assume a known factorization; instead, we diagnose whether existing state-based URL scores collapse observation-near states whose long-horizon futures differ because of local interaction variables.

\paragraph{Contact-rich dynamics and manipulation.}
Contact-rich manipulation is a canonical setting where local interaction variables change effective dynamics.
Classical work on pushing and manipulation emphasizes that contact mode, friction, and stability determine controllability and feasible motion~\citep{mason-ijrr86a,lynch-ijrr96a}.
Trajectory-optimization and contact-implicit methods explicitly model or optimize through contact transitions~\citep{posa-ijrr14a,manchester-isrr17a,patel-ral19a,aceituno-rss20a,jin-arxiv24a}.
Learning-based contact models also study the difficulty of representing discontinuous or hybrid contact dynamics~\citep{hogan-arxiv16a,pfrommer-corl20a}.
Our paper does not propose a contact model.
Instead, we use contact-rich cube manipulation as a controlled diagnostic for a broader URL question: whether learned future-reachability scores preserve interaction-conditioned successor distinctions.

\paragraph{Offline goal-conditioned benchmarks and diagnostics.}
Offline goal-conditioned RL evaluates whether policies or representations learned from static data can solve downstream goals without additional online collection~\citep{kostrikov-iclr22a,ma-neurips22a,park-iclr25a}.
Benchmarks such as \ogbench{} report aggregate goal-reaching success, which is useful for comparing algorithms but does not identify the physical stage at which a policy fails~\citep{park-iclr25a}.
Our diagnostic decomposes manipulation rollouts into approach, control acquisition, maintained control, and transport.
This separates coarse reachability from the local interaction transition and from the ability to preserve and exploit that interaction regime.